\chapter{Programs — Sqrt}
%%% generated by the blueprint pipeline from TCSlib.Programs.Light.Lib.Sqrt
\section{Overview}

Declarations of the Lean module \texttt{TCSlib.Programs.Light.Lib.Sqrt}.

\section{Declarations}

\begin{definition}[Argument variable index]
\lean{Light.Sqrt.Arg}
\label{Light.Sqrt.Arg}
\leanok
The argument variable is assigned index $0$.
\end{definition}

\begin{definition}[Candidate variable index]
\lean{Light.Sqrt.Cand}
\label{Light.Sqrt.Cand}
\leanok
\uses{Light.Sqrt.Arg}
The candidate variable is assigned index $1$.
\end{definition}

\begin{definition}[Square variable index]
\lean{Light.Sqrt.Square}
\label{Light.Sqrt.Square}
\leanok
\uses{Light.Sqrt.Arg}
The square variable is assigned index $2$.
\end{definition}

\begin{definition}[Integer square root program body]
\lean{Light.sqrtBody}
\label{Light.sqrtBody}
\leanok
\uses{Light.Sqrt.Arg, Light.Sqrt.Cand, Light.Sqrt.Square, Light.Stmt, Light.k, Light.v}
The program initializes the candidate variable to $0$ and the square variable to $1$.
While the square variable is at most the argument, it replaces the square by the square
plus twice the candidate plus $3$, then increases the candidate by $1$; when the loop
ends, it assigns the candidate to the argument variable.
\end{definition}

\begin{definition}[Square root running time bound]
\lean{Light.sqrtTime}
\label{Light.sqrtTime}
\leanok
For each natural number $K$, define the running time bound to be $18\lfloor\sqrt
K\rfloor+12$, where $\lfloor\sqrt K\rfloor$ is the integer square root of $K$.
\end{definition}

\begin{theorem}[Square root procedure correctness]
\lean{Light.sqrt_meets}
\label{Light.sqrt_meets}
\leanok
\uses{Light.Ends.next, Light.Ends.whileBlock, Light.Meets, Light.Meets.of_body, Light.frame, Light.sqrtBody, Light.sqrtTime, Light.update_frame_setLocal}
\difficulty{4}
Let a program contain the specified square root procedure, let $K$ be a natural number,
and let $\mu$ be an initial memory function. If the word limit is at least $3K+4$, then
on input $K$ the procedure runs within $18\lfloor\sqrt K\rfloor+12$ steps, returns
$\lfloor\sqrt K\rfloor$, and leaves memory unchanged.
\end{theorem}
